% Bagian: turing-machines
% Bab: machines-computations
% Subbagian: configuration

\documentclass[../../../include/open-logic-section]{subfiles}

\begin{document}

\olfileid[id]{cmp}{tur}{con}
\olsection{Konfigurasi dan Komputasi}

\begin{explain}
Ingatlah penelusuran konfigurasi-konfigurasi mesin genap sebelumnya.
Mekanisme khayalan yang terdiri atas pita, kepala baca-tulis, dan program mesin
Turing sebenarnya hanyalah cara intuitif untuk membayangkan komputasi mesin
Turing. Secara formal, kita dapat mendefinisikan komputasi suatu mesin Turing
pada masukan tertentu sebagai barisan \emph{konfigurasi}---dan pada gilirannya
suatu konfigurasi adalah barisan simbol (yang bersesuaian dengan isi pita pada
saat tertentu dalam komputasi), suatu bilangan yang menunjukkan posisi kepala
baca-tulis, dan suatu keadaan. Dengan menggunakan semua ini, kita dapat
mendefinisikan apa yang dihitung mesin Turing~$M$ pada masukan tertentu.
\end{explain}

\begin{defn}[Konfigurasi]
Suatu \emph{konfigurasi} mesin Turing $M = \tuple{Q, \Sigma, q_0,
\delta}$ adalah tripel $\tuple{C, m, q}$ dengan
\begin{enumerate}
\item $C \in \Sigma^*$ merupakan barisan berhingga simbol-simbol dari $\Sigma$,
\item $m \in \Nat$ merupakan bilangan $< \len{C}$, dan
\item $q \in Q$.
\end{enumerate}
Secara intuitif, barisan~$C$ adalah isi pita (simbol-simbol pada semua petak
mulai dari petak paling kiri hingga petak tak-kosong terakhir atau petak
terakhir yang pernah dikunjungi), $m$~adalah nomor petak yang sedang dipindai
kepala baca-tulis (dengan $0$ sebagai nomor petak paling kiri), dan $q$ adalah
keadaan mesin saat ini.
\end{defn}

\begin{explain}
Masukan yang mungkin bagi mesin Turing ialah suatu barisan simbol, biasanya
barisan yang mengodekan bilangan dalam bentuk tertentu. Konfigurasi awal mesin
Turing adalah konfigurasi tempat kita mulai menjalankan mesin Turing pada
masukan tersebut: pita memuat penanda ujung pita yang langsung diikuti masukan
yang tertulis pada petak-petak di sebelah kanan; kepala baca-tulis memindai
petak masukan paling kiri (yakni, petak di sebelah kanan penanda ujung kiri);
dan mekanisme berada dalam keadaan awal yang ditetapkan, yakni~$q_0$.
\end{explain}

\begin{defn}[Konfigurasi awal]
\emph{Konfigurasi awal} $M$ untuk masukan $I \in \Sigma^*$ adalah
\[
\tuple{\TMendtape \frown I \frown \TMblank, 1, q_0}.
\]
\end{defn}

\begin{explain}
Simbol~$\frown$ menyatakan \emph{konkatenasi}---untai masukan bermula tepat di
sebelah kanan penanda ujung kiri.
\end{explain}

\begin{defn}
Kita mengatakan bahwa suatu konfigurasi $\tuple{C, m, q}$ \emph{menghasilkan
konfigurasi $\tuple{C', m', q'}$ dalam satu langkah} (menurut~$M$) jika dan
hanya jika
\begin{enumerate}
\item simbol ke-$m$ dari $C$ adalah $\sigma$,
\item himpunan instruksi $M$ menentukan $\delta(q, \sigma) =
  \tuple{q', \sigma', D}$,
\item simbol ke-$m$ dari $C'$ adalah $\sigma'$, dan
\item
\begin{enumerate}
\item $D = L$ dan $m' = m - 1$ jika $m>0$, sedangkan jika tidak, $m'=0$; atau
\item $D = R$ dan $m' = m + 1$; atau
\item $D = N$ dan $m' = m$,
\end{enumerate}
\item jika $m' = \len{C}$, maka $\len{C'} = \len{C} + 1$ dan simbol ke-$m'$
  dari $C'$ adalah~$\TMblank$. Jika tidak, $\len{C'}=\len{C}$,
\item untuk setiap $i$ sedemikian sehingga $i < \len{C}$ dan $i \neq m$,
  berlaku $C'(i) = C(i)$.
\end{enumerate}
\end{defn}

\begin{defn}\ollabel{defn:run-output}
Suatu \emph{operasi $M$ pada masukan~$I$} adalah barisan berhingga atau tak
berhingga konfigurasi $C_0,C_1,\dots$ dari~$M$, dengan $C_0$ sebagai
konfigurasi awal $M$ bagi masukan~$I$, dan setiap $C_i$ menghasilkan
$C_{i+1}$ dalam satu langkah setiap kali $C_{i+1}$ ada dalam barisan tersebut.

Kita mengatakan bahwa $M$ \emph{berhenti pada masukan $I$ setelah $k$ langkah}
jika $C_k = \tuple{C, m, q}$, simbol ke-$m$ dari~$C$ adalah~$\sigma$, dan
$\delta(q, \sigma)$ tidak terdefinisi. Dalam hal itu, \emph{keluaran} $M$ bagi
masukan~$I$ adalah~$O$, dengan $O$ suatu untai simbol yang tidak berakhir
dengan~$\TMblank$ sedemikian sehingga
$C = \TMendtape \concat O \concat \TMblank^j$ untuk suatu~$j \in \Nat$.
($\TMblank^j$ adalah barisan yang terdiri atas $j$ simbol~$\TMblank$.)
\end{defn}

\begin{explain}
Menurut definisi ini, keluaran~$O$ dari~$M$ selalu berakhir dengan simbol selain
$\TMblank$; atau, jika pada saat~$k$ seluruh pita terisi~$\TMblank$ (kecuali
simbol~$\TMendtape$ paling kiri), $O$~adalah untai kosong.
\end{explain}

\end{document}
